\documentclass[../../main.tex]{subfiles}

\begin{document}

\section{弱序列紧性}\label{section:ch08-弱序列紧性}

\begin{theorem}\label{theorem:ch08-有界序列弱收敛子列}
令 $E$ 为可测集而 $1<p<\infty$. 则 $L^{p}(E)$ 中的每个有界序列有在 $L^{p}(E)$ 中弱收敛的子序列.
\end{theorem}

\begin{proof}
令 $\{f_{n}\}$ 为 $L^{p}(E)$ 中的有界序列. 设 $q$ 为 $p$ 的共轭. 我们必须找到一个子序列 $\{f_{n_{k}}\}$ 和函数 $f\in L^{p}(E)$, 使得对所有 $g\in L^{q}(E)$,
\begin{align*}
\lim_{k\to\infty}\int_{E}f_{n_{k}}\cdot g=\int_{E}f\cdot g.
\end{align*}
令 $\{h_{k}\}$ 为 $L^{q}(E)$ 的一个可数稠密子集 (见前一章定理 11). 由于 $\{f_{n}\}$ 有界, 对每个 $k$, 实数序列 $\{\int_{E}f_{n}\cdot h_{k}\}$ 有界. 根据对角论证法, 存在子序列 $\{f_{n_{j}}\}$ 使得对每个 $k$,
\begin{align}
\lim_{j\to\infty}\int_{E}f_{n_{j}}\cdot h_{k}\text{ 存在}.\label{align:ch08-对角极限}
\end{align}
定义在 $\{h_{k}\}$ 的线性扩张上的泛函 $T$ 为
\begin{align*}
T(h)=\lim_{j\to\infty}\int_{E}f_{n_{j}}\cdot h.
\end{align*}
该极限在所有 $h\in L^{q}(E)$ 上唯一延拓为有界线性泛函: 由于 $\{f_{n_{j}}\}$ 有界, 该泛函的有界性由 H\"older 不等式保证. 根据 Riesz 表示定理, 存在 $f\in L^{p}(E)$ 使得对所有 $h\in L^{q}(E)$, $T(h)=\int_{E}f\cdot h$. 于是 $\{f_{n_{j}}\}$ 弱收敛于 $f$.
\end{proof}

\begin{theorem}[Helly 定理]\label{theorem:ch08-Helly}
令 $1<p<\infty$. 则 $L^{p}(E)$ 的单位球是弱序列紧的: 单位球中的每个序列有在 $L^{p}(E)$ 中弱收敛的子序列, 且弱收敛的极限属于单位球.
\end{theorem}

\begin{proof}
前一个定理给出子序列的存在性. 弱收敛极限的范数不超过序列范数的上确界, 这由定理 7 的\eqref{align:ch08-弱收敛有界}保证.
\end{proof}

\begin{example}\label{example:ch08-L1无弱收敛子列}
$L^{1}(E)$ 中的有界序列不一定有弱收敛的子序列. 例如, 在 $[0,1]$ 上令 $f_{n}=n\chi_{(0,1/n)}$. 则 $\|f_{n}\|_{1}=1$ 对所有 $n$ 成立, 但 $\{f_{n}\}$ 没有弱收敛的子序列: 若有子序列弱收敛于 $f$, 则对 $g=\chi_{[0,1]}\in L^{\infty}[0,1]$, $\int_{0}^{1}f_{n_{k}}\cdot g=1\to\int_{0}^{1}f\cdot g$, 而对 $g=\chi_{[a,b]}$ ($a>0$), 当 $k$ 充分大时 $\int_{0}^{1}f_{n_{k}}\cdot g=0$, 因此 $\int_{a}^{b}f=0$ 对所有 $a>0$ 成立, 从而 $f=0$ a.e., 矛盾于 $\int_{0}^{1}f=1$.
\end{example}

\begin{exercise}\label{exercise:ch08-27}
证明: 若 $1<p<\infty$, 则 $L^{p}(E)$ 中的每个有界序列有弱收敛子序列.
\end{exercise}

\begin{exercise}\label{exercise:ch08-28}
令 $1<p<\infty$ 而 $\{f_{n}\}$ 为 $L^{p}(E)$ 中的有界序列. 证明: 若对所有 $g$ 在 $L^{q}(E)$ 的一个稠密子集中, $\lim_{n\to\infty}\int_{E}f_{n}\cdot g$ 存在, 则 $\{f_{n}\}$ 弱收敛.
\end{exercise}

\begin{exercise}\label{exercise:ch08-29}
举例说明 $L^{1}[0,1]$ 中的有界序列没有弱收敛的子序列.
\end{exercise}

\begin{exercise}\label{exercise:ch08-30}
令 $1<p<\infty$. 证明 $L^{p}(E)$ 的闭单位球是弱序列紧的.
\end{exercise}

\begin{exercise}\label{exercise:ch08-31}
证明 Helly 定理对 $p=1$ 不成立.
\end{exercise}

\begin{exercise}\label{exercise:ch08-32}
令 $1<p<\infty$. 证明 $L^{p}(E)$ 中的有界闭凸集是弱序列紧的.
\end{exercise}

\begin{exercise}\label{exercise:ch08-33}
举例说明 $L^{\infty}[0,1]$ 中的有界序列没有弱收敛的子序列.
\end{exercise}

\begin{exercise}\label{exercise:ch08-34}
令 $1<p<\infty$ 而 $\{f_{n}\}$ 为 $L^{p}(E)$ 中的有界序列. 证明: 若每个子序列有进一步的子序列弱收敛于同一个 $f$, 则整个序列弱收敛于 $f$.
\end{exercise}

\begin{exercise}\label{exercise:ch08-35}
令 $1<p<\infty$. 证明 $L^{p}(E)$ 中的有界序列 $\{f_{n}\}$ 弱收敛于 $f$ 当且仅当每个有收敛子列.
\end{exercise}

\begin{exercise}\label{exercise:ch08-36}
令 $1<p<\infty$ 而 $\{f_{n}\}$ 为 $L^{p}(E)$ 中的有界序列. 证明: 若对每个可测集 $A$, $\{\int_{A}f_{n}\}$ 收敛, 则 $\{f_{n}\}$ 弱收敛.
\end{exercise}

\begin{exercise}\label{exercise:ch08-37}
令 $1<p<\infty$. 证明: 若 $C$ 是 $L^{p}(E)$ 中的有界闭凸集, 则 $C$ 是弱序列紧的.
\end{exercise}

\end{document}